%%%%%%%%%%  scan idx 455  |  folio 442  %%%%%%%%%%
\SGAPrintPage{442}%
\textbf{Lemma.} \uline{Let} \(X'\xrightarrow{f}X\) \uline{be a projective and perfect morphism}, \uline{with} \(X\) \uline{quasi-compact}, \uline{and let} \(\underline{L}^\bullet\) \uline{be a perfect complex over} \(\underline{O}_{X'}\), \uline{acyclic in degrees} \(>0\). \uline{Then there exists an integer} \(n_0\) \uline{such that for} \(n\ge n_0\) \uline{the complex} \(\mathbb Rf_*\bigl(\underline{L}^\bullet\otimes_{\underline{O}_{X'}}\underline{O}_{X'}(n)\bigr)\) \uline{is acyclic in degrees} \(>0\).

Apply this lemma to the complex reduced to \(\underline{O}_{X'}\); there is therefore an integer \(n_0\) such that for \(n\ge n_0\), we have \(R^if_*(\underline{O}_{X'}(n))=0\) when \(i\ge1\). We then descend by induction to \(n=0\). For this, consider the exact sequence:
\[
0\longrightarrow
\underline{O}_{X'}(n+1)
\longrightarrow
\underline{O}_{X'}(n)
\longrightarrow
\underline{O}_{Y'}(n)
\longrightarrow0,
\]
and apply the functor \(f_*\) to it. We have \(\mathbb Rf_*(\underline{O}_{Y'}(n))=i_*\circ\mathbb Rg_*(\underline{O}_{Y'}(n))\). Since \(Y'=\bfP(\underline{N})\), \(\underline{O}_{Y'}(n)\) is \(g_*\)-acyclic for \(n\ge0\). Thus for every \(n\ge0\), and every \(i\ge1\):
\[
R^if_*(\underline{O}_{Y'}(n))=0,
\]
and consequently we have the exact sequences:
\[
R^if_*(\underline{O}_{X'}(n+1))
\longrightarrow
R^if_*(\underline{O}_{X'}(n))
\longrightarrow0,
\]
which complete the induction.

ii) We must show that the canonical morphism:
\[
\bigoplus_{n\ge0}\underline{J}^n
\longrightarrow
\bigoplus_{n\ge0}f_*\bigl(\underline{O}_{X'}(n)\bigr)
\]
is an isomorphism. Since the property is local, we may suppose that \(\underline{J}\) is defined by a sequence \(f_1,\ldots,f_d\) of sections of \(\underline{O}_X\) such that the corresponding homomorphism \((\underline{O}_X)^d\longrightarrow\underline{O}_X\) is regular. Put: \(\underline{S}=\bigoplus_{n\ge0}\underline{J}^n\); we define a surjective algebra homomorphism: \(\underline{O}_X[T_1,\ldots,T_d]\xrightarrow{v}\underline{S}\), by sending

%%%%%%%%%%  scan idx 456  |  folio 443  %%%%%%%%%%
\SGAPrintPage{443}%
\(T_i\) to \(f_i^{(1)}\). This homomorphism defines a closed immersion \(X'\longrightarrow\bfP_X^{d-1}\), regular by 1.8 ii).

By 1.8 iii), \(\underline{S}\) is perfect as a graded \(\underline{O}_X[T_1,\ldots,T_d]\)-Module. In particular, \(\underline{S}\) is finitely presented as a graded Module. Let therefore:
\[
\underline{L}''\longrightarrow \underline{L}'\longrightarrow \underline{S}\longrightarrow0
\]
be a finite graded presentation of \(\underline{S}\) over \(\underline{O}_X[T_1,\ldots,T_d]\). On \(\bfP_X^{d-1}\) we obtain the exact sequence:
\[
\underline{M}''\longrightarrow \underline{M}'\longrightarrow\underline{O}_{X'}\longrightarrow0,
\]
where \(\underline{M}''=\underline{\Proj}(\underline{L}'')\), \(\underline{M}'=\underline{\Proj}(\underline{L}')\), so these are locally free of finite type over \(\bfP_X^{d-1}\). We can therefore write, for every \(n\), the commutative diagram:
\[
\begin{tikzcd}[column sep=large,row sep=large]
\underline{L}''_n \arrow[r] \arrow[d,"\sim"'] &
\underline{L}'_n \arrow[r] \arrow[d,"\sim"'] &
\underline{J}^n \arrow[r] \arrow[d] &
0\\
p_*(\underline{M}''(n)) \arrow[r] &
p_*(\underline{M}'(n)) \arrow[r] &
p_*(\underline{O}_{X'}(n)) \arrow[r] &
0,
\end{tikzcd}
\]
where \(p\) is the structural morphism: \(\bfP_X^{d-1}\to X\), and \(\underline{L}''_n\), \(\underline{L}'_n\) are the degree \(n\) components of \(\underline{L}''\), \(\underline{L}'\). The upper row is exact by hypothesis; the first two vertical arrows are isomorphisms by the classical results for the direct images of \(\underline{O}_P(n)\), \(P\) being a projective bundle (EGA III, 2.1.15); finally, the last row is exact when \(n\) is sufficiently large, since the kernel of \(\underline{M}'\to\underline{O}_{X'}\) is a perfect \(\underline{O}_P\)-Module, and we can apply the lemma recalled above, which shows that this kernel, tensored with \(\underline{O}_P(n)\) for \(n\) sufficiently large, is \(p_*\)-acyclic; likewise for the kernel of \(\underline{M}''\to \underline{M}'\). It follows that for \(n\) sufficiently large, \(\underline{J}^n\longrightarrow p_*(\underline{O}_{X'}(n))\) is an isomorphism. Since \(p_*(\underline{O}_{X'}(n))=f_*(\underline{O}_{X'}(n))\), we therefore have

%%%%%%%%%%  scan idx 457  |  folio 444  %%%%%%%%%%
\SGAPrintPage{444}%
the result for \(n\ge n_0\).

It remains to descend to \(n=0\). We proceed by descending induction, using the commutative diagram:
\[
\begin{tikzcd}[column sep=large,row sep=large]
0 \arrow[r] \arrow[d] &
\underline{J}^{n+1} \arrow[r] \arrow[d,"\sim"'] &
\underline{J}^n \arrow[r] \arrow[d] &
\underline{J}^n/\underline{J}^{n+1} \arrow[r] \arrow[d,"\sim"'] &
0 \arrow[d]\\
0 \arrow[r] &
f_*(\underline{O}_{X'}(n+1)) \arrow[r] &
f_*(\underline{O}_{X'}(n)) \arrow[r] &
f_*(\underline{O}_{Y'}(n)) \arrow[r] &
0.
\end{tikzcd}
\]
The bottom row is exact because \(R^1f_*(\underline{O}_{X'}(n+1))=0\) for every \(n\ge0\), by i), already proved. Furthermore, we have:
\[
f_*(\underline{O}_{Y'}(n))
=i_*\bigl(g_*(\underline{O}_{Y'}(n))\bigr)
=i_*(\mathbf S(\underline{N})_n)
=\underline{J}^n/\underline{J}^{n+1}.
\]
The arrow on the right is therefore an isomorphism; the one on the left is an isomorphism by the induction hypothesis; hence the middle one is also an isomorphism.

\sgaStableTarget{sga6.e07.prop.3-6}\textbf{Proposition 3.6.} \uline{With the hypotheses and notation of} 3.1 \uline{and} 3.3, \uline{we have the relation}:
\sgaStableTarget{sga6.e07.eq.3-6-1}\[
f_*\circ f^*=\Id_{K^{\bullet}(X)}.
\]

Let \(x\in K^{\bullet}(X)\); by the projection formula, we can write:
\[
f_*(f^*(x))=f_*(1.f^*(x))=x.f_*(1).
\]
It therefore suffices to show that \(f_*(1)=1\).

Now 3.5 shows that the cohomology of \(\mathbb Rf_*(\underline{O}_{X'})\) vanishes in degrees \(>0\), and equals \(\underline{O}_X\) in degree \(0\). The proposition follows at once.

\sgaStableTarget{sga6.e07.thm.3-7}\textbf{Theorem 3.7.} \uline{Let} \(i:Y\to X\) \uline{be a regular closed immersion}, \(X'\) \uline{the scheme obtained by blowing up} \(Y\), \uline{and} \(Y'=X'\times_XY\), \uline{so that we have the commutative diagram}

%%%%%%%%%%  scan idx 458  |  folio 445  %%%%%%%%%%
\SGAPrintPage{445}%
\[
\begin{tikzcd}[column sep=large,row sep=large]
Y' \arrow[r,"j"] \arrow[d,"g"'] & X' \arrow[d,"f"]\\
Y \arrow[r,"i"] & X.
\end{tikzcd}
\]
\uline{Suppose that} \(X\) \uline{is quasi-compact and has an ample sheaf}. \uline{Then}, \uline{we have the exact sequence}:
\sgaStableTarget{sga6.e07.eq.3-7-0}\[
0\longrightarrow K^{\bullet}(Y)\xrightarrow{u}K^{\bullet}(Y')\times K^{\bullet}(X)
\xrightarrow{v}K^{\bullet}(X'),
\]
\uline{where we put}, \uline{for} \(y\in K^{\bullet}(Y)\):
\sgaStableTarget{sga6.e07.eq.3-7-u}\[
u(y)=\bigl(g^*(y).\lambda_{-1}(F),-i_*(y)\bigr),
\]
\uline{and for} \(y'\in K^{\bullet}(Y')\), \(x\in K^{\bullet}(X)\):
\sgaStableTarget{sga6.e07.eq.3-7-v}\[
v(y',x)=j_*(y')+f^*(x),
\]
\uline{with the notation of} 3.1 \uline{and} 3.3. \uline{Moreover}, \(u\) \uline{has a left inverse} \(u'\), \uline{defined by} \(u'(y',x)=g_*(y')\).

\uline{If in addition} \(X\) \uline{is noetherian}, \uline{and} \(Y\) \uline{regular}, \uline{then} \(v\) \uline{is surjective}, \uline{so we have the exact sequence}:
\sgaStableTarget{sga6.e07.eq.3-7-2}\[
0\longrightarrow K^{\bullet}(Y)\xrightarrow{u}K^{\bullet}(Y')\times K^{\bullet}(X)
\xrightarrow{v}K^{\bullet}(X')\longrightarrow0.
\]

First, \(u\) has the left inverse \(u'\), and is a fortiori injective. Indeed, we have:
\[
g_*\bigl(g^*(y).\lambda_{-1}(F)\bigr)
=y.g_*(\lambda_{-1}(F))
\]
by the projection formula. Furthermore, \(g_*(\lambda_{-1}(F))=1\) (VI 5.9), which gives the result.

Next, \(v\circ u=0\). Indeed, for \(y\in K^{\bullet}(Y)\) we have:
\[
v\circ u(y)
=v\bigl(g^*(y)\lambda_{-1}(F),-i_*(y)\bigr)
=j_*\bigl(g^*(y)\lambda_{-1}(F)\bigr)-f^*(i_*(y))=0,
\]
by 3.4.

It therefore remains to show that \(\Ker(v)\subset \Im(u)\). Let therefore \(x\in K^{\bullet}(X)\), \(y'\in K^{\bullet}(Y')\), be such that:

%%%%%%%%%%  scan idx 459  |  folio 446  %%%%%%%%%%
\SGAPrintPage{446}%
\[
j_*(y')+f^*(x)=0.
\]
We must find \(y\in K^{\bullet}(Y)\) such that:
\sgaStableTarget{sga6.e07.eq.3-7-1}\begin{equation}\tag{3.7.1}
x=-i_*(y);\qquad y'=g^*(y).\lambda_{-1}(F).
\end{equation}
Since \(u\) has a left inverse \(u'\), we necessarily have
\[
y=u'(y',x)=g_*(y').
\]
Thus we must show that \(y\) satisfies the relations (3.7.1).

To show that \(x=-i_*(y)\), start with the relation \(f^*(x)=-j_*(y')\) and apply \(f_*\) to it. We therefore obtain: \(f_*(f^*(x))=-f_*(j_*(y'))=-i_*(g_*(y'))=-i_*(y)\) by the definition of \(y\). By Proposition 3.6, however, \(f_*(f^*(x))=x\); hence the required relation.

To complete the proof, we must show that
\[
y'=g^*(g_*(y'))\lambda_{-1}(F).
\]
Put \(z=y'-g^*(g_*(y'))\lambda_{-1}(F)\). We obtain: \(g_*(z)=g_*(y')-g_*(y')\,g_*(\lambda_{-1}(F))=0\). On the other hand, by hypothesis we have: \(j_*(y')=-f^*(x)\), whence, by the preceding calculation, \(j_*(y')=f^*(i_*(y))\), and by Proposition 3.4: \(j_*(y')=j_*\bigl(g^*(y)\lambda_{-1}(F)\bigr)\); consequently, we also have \(j_*(z)=0\). Now \(\Ker(j_*)\cap\Ker(g_*)=0\). Indeed, if \(j_*(z)=0\), then \(j^*(j_*(z))=0\); but \(j^*(j_*(z))=z(1-L)\) (Prop. 2.7), and for \(z(1-L)=0\) to hold there must exist \(z'\in K^{\bullet}(Y)\) such that \(z=z'\lambda_{-1}(F)\) (VI,5.10); applying \(g_*\) gives \(z'=g_*(z)=0\), hence \(z=0\).

Suppose now that \(X\) is noetherian and \(Y\) regular. Then \(Y'\) is also regular, since it is a projective bundle over \(Y\), and for every point \(y'\) of \(Y'\), the local ring \(\underline{O}_{X',y'}\) of \(X'\) at \(y'\) is regular, since \(\underline{O}_{Y',y'}\) is the quotient of \(\underline{O}_{X',y'}\) by a regular element (see EGA \(0_{\mathrm{IV}}\) 17.1.8).
